\documentclass[10pt,a4paper]{article}
\usepackage[utf8]{inputenc}
\usepackage[T1]{fontenc}
\usepackage[margin=1.8cm,top=2.6cm,bottom=2.2cm]{geometry}
\usepackage{parskip}
\usepackage{microtype}
\usepackage{amsmath,amssymb,mathtools}
\usepackage{array,booktabs,tabularx,multirow}
\usepackage[table]{xcolor}
\usepackage{graphicx,float}
\usepackage{titlesec}
\usepackage{enumitem}
\usepackage{fancyhdr}

% ---- Section numbering: Roman numerals ---------------------------------
\renewcommand\thesection{\Roman{section}}
% ---- Equation numbering: (I.1), (II.2), etc. ---------------------------
\numberwithin{equation}{section}

% ---- Section formatting ------------------------------------------------
\titleformat{\section}{\large\bfseries}{\thesection.\;}{0em}{}[\titlerule]
\titleformat{\subsection}{\normalsize\bfseries}{}{0em}{}
\titleformat{\subsubsection}{\normalsize\itshape}{}{0em}{}
\titlespacing{\section}{0pt}{1.6ex plus .2ex}{0.6ex}
\titlespacing{\subsection}{0pt}{1.0ex}{0.3ex}

% ---- Page headers ------------------------------------------------------
\pagestyle{fancy}
\fancyhf{}
\setlength{\headheight}{15pt}
\fancyhead[L]{\small\textit{Layer~7: SAM Chemistry}}
\fancyhead[R]{\small Aaran Nihalani $\bullet$ Eton iGEM Dry Lab}
\fancyfoot[C]{\small\thepage}
\renewcommand{\headrulewidth}{0.35pt}

% ---- Macros ------------------------------------------------------------
\newcommand{\cit}[1]{[#1]}   % inline citation: [1], [7], etc.
\newcommand{\Mms}{\text{M}^{-1}\text{s}^{-1}}
\newcommand{\As}{\text{\AA}^{-1}}

\begin{document}

%=======================================================================
\begin{center}
{\LARGE\bfseries Layer 7: SAM Chemistry}\\[0.4em]
{\large Three-Component Probe-Thiol/C6FC/MCH Monolayer: Deposition, Stability, and Ca$^{2+}$ Artefact}\\[0.25em]
{\small\color{gray} Eton iGEM 2026, Dry Lab \quad$|$\quad
Layer 7 of 17-Layer Biosensor Pipeline \quad$|$\quad
\texttt{sam\_loading.py}, \texttt{sam\_stability.py}, \texttt{calcium\_artefact.py}}
\end{center}
\vspace{0.4em}\hrule\vspace{0.7em}

\noindent
Layer~7 characterises the self-assembled monolayer (SAM) that immobilises the peptide
probe on the gold electrode surface. Its primary output is the probe surface density
$\Gamma_p$, which passes to Layer~8 (steric accessibility) and Layer~9 (surface kinetics).
Secondary outputs are bounds on SAM degradation over the 60-minute assay window and a
quantitative account of the Ca$^{2+}$ artefact and its ratiometric cancellation.

Every parameter value must be confirmed from a paper physically present in
\texttt{papers/layer\_07/}; values without a paper source are explicitly flagged throughout.

%=======================================================================
\section{Mixed SAM Architecture}
\label{sec:architecture}

\subsection{Molecular components}

The sensor electrode carries a three-component thiol monolayer. The probe-thiol consists
of a peptide substrate bearing a methylene-blue (MB) redox label, tethered to gold via
a C$_6$ alkyl spacer and a PEG$_6$ flexible linker
(HS--(CH$_2$)$_6$--PEG$_6$--SGKIPRTLTA--MB). MMP-9 (or HNE on the secondary
channel) cleaves the peptide, releasing the MB-bearing fragment into solution and
producing a signal-off decrease in SWV peak current proportional to enzymatic activity
(Layer~0, Section~V).

A co-immobilised reference reporter, 6-ferrocenylhexanethiol (C6FC,
Fc--(CH$_2$)$_6$--SH), provides an internal electrochemical standard. C6FC shares the
same C$_6$ alkyl chain as the probe-thiol spacer, ensuring compatible packing geometry
in the SAM. The ferrocene (Fc) headgroup gives a reversible one-electron redox signal.
Adj\'{e}mian et~al.\ (2010)~\cit{12} report a formal potential of
$\sim$150--155~mV vs.\ SCE ($\approx +0.20$~V vs.\ Ag/AgCl) for an alkyl-ferrocene
label on an analogous Fc-peptide-PEG monolayer on gold (p.~10350 of \cit{12});
this value is expected to be close to the C6FC formal potential but must be confirmed
directly on the fabricated device. The Fc signal is well positive of the MB reduction
peak; the two peaks are resolved in a single SWV scan without additional switching
(Section~\ref{sec:c6fc}). The MB formal potential in the assay buffer is not confirmed
from any paper in \texttt{papers/layer\_07/} and must be determined by cyclic
voltammetry on the fabricated electrode.

The diluent is 6-mercapto-1-hexanol (MCH, HS--(CH$_2$)$_6$--OH). MCH fills gold sites
not occupied by the probe-thiol or C6FC, presenting a chemically inert surface to the
wound fluid and preventing non-specific adsorption of wound proteins to the bare gold.

\subsection{Close-packed monolayer density}

The maximum surface density of a fully close-packed alkanethiol monolayer on gold,
measured by reductive desorption coulometry, is \cit{6}:

\begin{equation}
  \Gamma_{\max} = 9.3 \times 10^{-10}\ \text{mol~cm}^{-2} = 930\ \text{pmol~cm}^{-2}
  \label{eq:gamma_max}
\end{equation}

Love et al.\ report an independent cross-check of approximately
750~pmol~cm$^{-2}$ for densely packed alkanethiol monolayers on Au(111) \cit{1}.
Equation~\eqref{eq:gamma_max} is an upper bound. The probe-thiol cannot approach
$\Gamma_{\max}$ due to the bulky peptide headgroup; the target probe coverage
$\Gamma_p = 2$--5~pmol~cm$^{-2}$ is 0.2--0.5\% of the close-packed limit.

The following assumptions underlie this section:

\begin{enumerate}[noitemsep]
  \item The probe-thiol adsorbs to gold through the thiol--gold bond; MCH backfilling
    fills the remaining gold surface.
  \item $\Gamma_{\max} = 930$~pmol~cm$^{-2}$ is an upper bound on MCH packing; the
    actual MCH density is lower because probe-thiols occupy a fraction of the surface.
\end{enumerate}

\subsection{C6FC reference reporter and ratiometric signal}
\label{sec:c6fc}

C6FC (Fc--(CH$_2$)$_6$--SH) is co-immobilised with the probe-thiol at a molar ratio
that is to be optimised on the fabricated electrode. Its function is threefold.

\begin{enumerate}[noitemsep]
  \item \textbf{Internal area reference.} The Fc signal $I_{\text{Fc}}$ is proportional
    to the number of C6FC molecules on the electrode and is insensitive to MMP-9 activity,
    because C6FC contains no peptide substrate. $I_{\text{Fc}}$ therefore tracks
    electrode area and surface loading, providing a normalisation factor that cancels
    run-to-run variation in electrode geometry or deposition yield.
  \item \textbf{Fouling correction.} Non-specific protein adsorption attenuates both
    the MB and Fc signals by approximately the same factor (fouling layer acts as a
    diffusion barrier). Taking the ratio $I_{\text{MB}}/I_{\text{Fc}}$ suppresses this
    common-mode attenuation.
  \item \textbf{Ratiometric readout.} The device signal is the ratio:
    \begin{equation}
      R(t) = \frac{I_{\text{MB}}(t)}{I_{\text{Fc}}}
      \label{eq:ratio}
    \end{equation}
    where $I_{\text{Fc}}$ is measured simultaneously with $I_{\text{MB}}(t)$ in the
    same SWV scan. The normalised fractional cleavage is:
    \begin{equation}
      \frac{\Delta R}{R_0} = \frac{R_0 - R(t)}{R_0} = f_{\text{cleaved}}(t)
      \label{eq:ratio_signal}
    \end{equation}
    provided $I_{\text{Fc}}$ is stable over the assay window (see assumption below).
\end{enumerate}

The electron-transfer kinetics of Fc-terminated alkanethiol chains on gold are
characterised in Smalley et~al.\ (2003)~\cit{8}: the tunnelling decay constant through
the alkyl bridge is $\beta = 0.94 \pm 0.06\ \text{\AA}^{-1}$, confirming that ET from
Fc to gold through the C$_6$ chain is fast relative to the SWV timescale and produces
a sharp, well-defined peak. The same $\beta$ value applies to the MB probe-thiol chain
(Section~\ref{sec:calcium}).

The following assumption underlies this subsection. MMP-9 cleaves SGKIPRTLTA but has no
known activity against the Fc--(CH$_2$)$_6$--SH molecule, which carries no peptide
substrate; $I_{\text{Fc}}$ is therefore expected to be stable over the 60-minute assay
window. This has not been confirmed by a control experiment on the fabricated electrode
and should be verified before ratiometric calibration is finalised. The target C6FC
surface density relative to probe-thiol density is not confirmed from any paper in
\texttt{papers/layer\_07/} and must be optimised empirically.

%=======================================================================
\section{Deposition Protocol and Target Surface Density}
\label{sec:deposition}

\subsection{Confirmed deposition protocol}

Wcis{\l}o et~al.\ (2021)~\cit{11} describe immobilisation of a cysteine-terminated
peptide substrate (SH--(CH$_2$)$_{10}$--CO--Arg-Val-Arg-Arg) on gold by the
two-step procedure applicable to this device:

\begin{enumerate}[noitemsep]
  \item Immerse polished gold electrode in 5~mM degassed ethanolic solution of the
    peptide-thiol for 20~h at room temperature.
  \item Rinse with ethanol and deionised water; dry under N$_2$.
  \item Backfill with 2~mM ethanolic 1-octanethiol (analogous to MCH) for 60~min,
    to eliminate pinholes and block non-specific adsorption.
  \item Optional protein block (BSA, 1\% in PBS, 60~min at 4$^\circ$C) if wound-matrix
    fouling is anticipated.
\end{enumerate}

For this device, 1-octanethiol is replaced by MCH (6-mercapto-1-hexanol, HS--(CH$_2$)$_6$--OH)
to match the C$_6$ alkyl chain length of the probe-thiol spacer. The target deposition
concentration (5~mM) and time (20~h) are adopted from Wcis{\l}o~2021~\cit{11}; exact
optimisation for the specific SGKIPRTLTA--MB probe must be confirmed by measuring
$\Gamma_p$ on the fabricated electrode.

\subsection{Langmuir adsorption framework}
\label{sec:langmuir}

The equilibrium coverage of a thiol on gold from a solution of molar concentration $c$
is described by the Langmuir isotherm:

\begin{equation}
  \Gamma_p = \Gamma_{\max} \cdot \frac{K_L\,c}{1 + K_L\,c}
  \label{eq:langmuir}
\end{equation}

where $K_L = k_{\text{ads}}/k_{\text{des}}$ (M$^{-1}$) is the Langmuir adsorption
constant and $\Gamma_{\max}$ is the monolayer saturation coverage.

\textbf{Saturation coverage.} For the directly analogous Fc-tetrapeptide-PEG$_7$ system
on gold (PEG spacer, short peptide sequence, Fc label, MCH backfill), Adj\'{e}mian~2010
\cit{12} report saturation at $\Gamma_{\max} \approx 6 \times 10^{-11}\ \text{mol~cm}^{-2}
= 60\ \text{pmol~cm}^{-2}$ (Experimental Section, p.~10349 of \cit{12}). This is
substantially lower than the alkanethiol limit ($930\ \text{pmol~cm}^{-2}$,
Eq.~\eqref{eq:gamma_max}) because the bulky PEG--peptide headgroup sterically limits
packing. This value is used as an estimate of $\Gamma_{\max}$ for the SGKIPRTLTA--MB
probe, pending direct coulometric measurement on the fabricated electrode.

\textbf{Effectively irreversible adsorption.} Thiol adsorption on gold is effectively
irreversible under typical deposition conditions: the desorption rate constant
$k_{\text{des}}$ is negligible on the timescale of SAM preparation, giving
$K_L \gg 1/c$ at any practical concentration. The Langmuir isotherm
(Eq.~\eqref{eq:langmuir}) therefore sits on its saturating plateau
($\Gamma_p \approx \Gamma_{\max}$) at both the micromolar concentrations used in
Adj\'{e}mian~2010 \cit{12} and the 5~mM concentration adopted from Wcis{\l}o~2021
\cit{11}. Coverage is consequently governed by \emph{kinetic} Langmuir rather than
equilibrium concentration:

\begin{equation}
  \frac{d\Gamma_p}{dt} = k_{\text{ads}}\,c\,(\Gamma_{\max} - \Gamma_p)
  \label{eq:langmuir_kin}
\end{equation}

The solution is $\Gamma_p(t) = \Gamma_{\max}(1 - e^{-k_{\text{ads}}ct})$: coverage
increases toward $\Gamma_{\max}$ at a rate set by exposure time, not by the position
of the equilibrium isotherm. This is precisely the regime exploited by Adj\'{e}mian~2010,
who achieve controlled sub-saturating coverages
($\Gamma_p \approx 0.4\text{--}2 \times 10^{-11}\ \text{mol~cm}^{-2}$) by limiting
the immersion time to 30--120~s at 2--5~µM \cit{12}, rather than by adjusting
concentration. The Wcis{\l}o 20~h protocol at 5~mM drives the system fully to
$\Gamma_{\max}$; working coverage well below saturation must therefore be achieved by
controlling immersion time, not by relying on equilibrium partitioning.

\textbf{Working coverage selection.} The device requires
$\Gamma_p = 2\text{--}5\ \text{pmol~cm}^{-2}$ (enzyme-access constraint,
Section~\ref{sec:deposition}), substantially below $\Gamma_{\max} = 60\ \text{pmol~cm}^{-2}$.
Equation~\eqref{eq:langmuir_kin} implies this is reached at short exposure times
(or very dilute probe concentration), consistent with the Adj\'{e}mian~2010 approach.
The exact immersion time that yields $\Gamma_p = 2\text{--}5\ \text{pmol~cm}^{-2}$
for the SGKIPRTLTA--MB probe must be calibrated on the fabricated electrode by
coulometric integration of the reductive desorption charge \cit{6}.

\textbf{Three-component competitive Langmuir.} With probe-thiol, C6FC, and MCH all
present in the deposition solution, the three species compete for gold sites.
The competitive Langmuir isotherm for component $i$ is:

\begin{equation}
  \theta_i = \frac{K_{L,i}\,c_i}{1 + K_{L,\text{probe}}\,c_{\text{probe}}
                                    + K_{L,\text{C6FC}}\,c_{\text{C6FC}}
                                    + K_{L,\text{MCH}}\,c_{\text{MCH}}}
  \label{eq:langmuir_comp}
\end{equation}

The coverage ratio $\Gamma_{\text{probe}}/\Gamma_{\text{C6FC}}$ therefore depends on
$K_{L,\text{probe}}\,c_{\text{probe}}\,/\,K_{L,\text{C6FC}}\,c_{\text{C6FC}}$.
The individual Langmuir constants $K_{L,\text{probe}}$, $K_{L,\text{C6FC}}$, and
$K_{L,\text{MCH}}$ for the specific thiols used here are not confirmed from any paper in
\texttt{papers/layer\_07/}; the deposition solution molar ratios must therefore be
optimised empirically. Once deposited, the individual surface densities
$\Gamma_{\text{probe}}$ and $\Gamma_{\text{C6FC}}$ should be verified by reductive
desorption coulometry (Widrig~\cit{6}) or by integrating the separate CV peaks of MB
and Fc in the same post-deposition scan.

\subsection{Confirmed functional coverage limits}

Adj\'{e}mian et~al.\ (2010)~\cit{12} characterised the relationship between surface
coverage $\Gamma_0$ and cleavage efficiency for a ferrocene (Fc)-labelled tetrapeptide
end-grafted on gold via a PEG$_7$ linker, using MCH backfill (architecture directly
analogous to the MB-SGKIPRTLTA--PEG$_6$ probe). Their key finding (Fig.~2 and text):

\begin{itemize}[noitemsep]
  \item At saturated undiluted layers ($\Gamma_0 \approx 6 \times 10^{-11}$~mol~cm$^{-2}
    = 60\ \text{pmol~cm}^{-2}$), cleavage yield drops to $\sim$50\% because chains form
    compact domains inaccessible to the enzyme.
  \item For $\Gamma_0 \leq 10^{-11}$~mol~cm$^{-2} = 10\ \text{pmol~cm}^{-2}$ with MCH
    backfill, cleavage yield exceeds 80\%.
  \item The working coverage used throughout Adj\'{e}mian~2010 is
    $(0.5\text{--}1) \times 10^{-11}$~mol~cm$^{-2} = 5\text{--}10\ \text{pmol~cm}^{-2}$.
\end{itemize}

The device target $\Gamma_p = 2$--5~pmol~cm$^{-2}$ therefore sits within the confirmed
functional regime. Adj\'{e}mian~2010 also confirms that surface kinetics follow the
Michaelis--Menten two-step mechanism (Layer~9), with enzymatic constants measurable
directly from the electrode response.

\begin{table}[H]
\centering
\small\renewcommand{\arraystretch}{1.3}
\begin{tabular}{@{}p{4.5cm}p{5.0cm}p{4.0cm}@{}}
\toprule
\textbf{Parameter} & \textbf{Status} & \textbf{Source} \\
\midrule
Deposition protocol & Confirmed (5~mM EtOH, 20~h, MCH backfill) & Wcis\l{}o~2021~\cit{11} \\
Functional $\Gamma_p$ limit & Confirmed ($\leq 10\ \text{pmol~cm}^{-2}$) & Adj\'{e}mian~2010~\cit{12} \\
$k_{\text{cat}}$, $K_m$ for SGKIPRTLTA & Confirmed (86.4~s$^{-1}$, 1.3~mM) & Kridel~2001~\cit{10} \\
Absolute $\Gamma_p$ on fabricated electrode & \textbf{Unconfirmed}: must be measured & Coulometric integration \\
Surface $k_{\text{eff,surf}}$ & \textbf{Unconfirmed}: must be measured on device & Layer~9 \\
\bottomrule
\end{tabular}
\caption{Status of probe-loading parameters. The two unconfirmed items require measurements
  on the fabricated electrode; they cannot be borrowed from solution or analogous probes.}
\label{tab:params_status}
\end{table}

\subsection{Probe spacing constraint from enzyme size}

Once $\Gamma_p$ is established from the literature, the average probe spacing is:

\begin{align}
  \ell &= (\Gamma_p \cdot N_A)^{-1/2} \label{eq:spacing}
\end{align}

For the sensor to function, MMP-9 (molecular weight 92~kDa, hydrodynamic radius
$\approx 3$--4~nm) must approach and bind each surface-immobilised peptide. This
imposes the requirement $\ell \gtrsim 2 \times r_{\text{MMP-9}} \approx 6$--8~nm,
corresponding to $\Gamma_p \lesssim 3$~pmol~cm$^{-2}$. At higher probe densities,
neighbouring probe molecules sterically block enzyme access, reducing the effective
cleavage rate and underestimating free MMP-9 activity. This constraint must be confirmed
with a cleavage-efficiency vs.\ coverage measurement on the fabricated electrode.

The following assumptions underlie this section:

\begin{enumerate}[noitemsep]
  \item Probe-thiol adsorbs to gold via the thiol--gold bond; MCH backfilling fills
    remaining surface sites.
  \item A target of $\Gamma_p \lesssim 3$~pmol~cm$^{-2}$ is adopted on steric grounds
    (enzyme-access requirement); this value is not confirmed from any paper in
    \texttt{papers/layer\_07/} and is flagged as Limitation~1.
\end{enumerate}

%=======================================================================
\section{SAM Stability over the Assay Window}
\label{sec:stability}

Three mechanisms can degrade the SAM during the 60-minute assay: oxidative desorption
by H$_2$O$_2$, thiol exchange with free thiols in the wound matrix, and electrochemical
desorption by the SWV scan potential.

\textbf{Applicability of alkanethiol literature.} All three subsections draw on papers
that measured simple $n$-alkanethiol SAMs: hydrophobic chains with no polar headgroups,
minimal solvent exposure, and compact crystalline packing on gold. The probe-thiol in
this device is not a simple alkanethiol: it carries a PEG$_6$ polyether linker and a
10-residue charged peptide headgroup, both hydrophilic and solvent-swollen at pH~7.4.
The MCH diluent (HS--(CH$_2$)$_6$--OH, $n = 6$) is structurally closer to the literature
systems but is interdigitated with the bulky probe-thiol at 2--5~pmol~cm$^{-2}$, which
disrupts the close-packed alkanethiol geometry on which published rate constants and
desorption potentials were measured. The alkanethiol data below provide
order-of-magnitude bounds only and are not quantitatively applicable to this mixed
peptide-thiol/MCH monolayer without experimental confirmation on the fabricated electrode.

\subsection{Oxidative desorption}

Schoenfisch and Pemberton characterised alkanethiol SAM oxidation under ambient conditions
\cit{5}. Their central finding is that ozone (O$_3$), not O$_2$ or UV, is the primary
oxidant in air, converting Au--SR to Au$^0$ + RSO$_x$H. Under ambient O$_3$ at 28--60~ppb,
a C$_{12}$ SAM on gold lost 36.9\% of its total sulphur in 1~h (Table~5 of \cit{5}).

In the present device, the oxidant is H$_2$O$_2$ at 10--100~$\mu$M in aqueous wound fluid;
ozone is absent from the aqueous phase. The relative oxidising power of the two species
is reflected in their standard reduction potentials:

\begin{equation}
  E^\circ(\text{O}_3/\text{O}_2) = +2.07\ \text{V},
  \qquad
  E^\circ(\text{H}_2\text{O}_2/\text{H}_2\text{O}) = +1.76\ \text{V}
  \label{eq:redox}
\end{equation}

The 310~mV thermodynamic advantage of O$_3$ corresponds to a rate advantage of several
orders of magnitude for thiolate oxidation. The Schoenfisch data therefore provide a
conservative upper bound: even under the harshest oxidative conditions tested (ambient O$_3$
in air), SAM loss is approximately 37\% in 1~h for a C$_{12}$ chain. Aqueous H$_2$O$_2$
at wound-fluid concentrations is a substantially weaker oxidant; oxidative SAM loss over
the 60-minute assay is not expected to approach this bound.

No quantitative rate constant $k_{\mathrm{ox}}$ for H$_2$O$_2$ oxidation of thiol SAMs
in solution is available in \texttt{papers/layer\_07/}. The argument above is comparative
only; see Limitation~2.

\subsection{Thiol exchange}

Schlenoff et al.\ investigated the exchange of adsorbed alkanethiols on gold by free
thiol in solution \cit{7}. The key finding is that the exchange rate shows ``very little
dependence of rate on thiol concentration,'' meaning the process is not bimolecular but
involves a rate-limiting surface step. In hexane, the pseudo-first-order exchange rate
constant is \cit{7}:

\begin{equation}
  k_1 = 1.7 \times 10^{-5}\ \text{s}^{-1} \quad (\text{hexane,}\ 22^\circ\text{C})
  \label{eq:k1}
\end{equation}

The fraction exchanged at time $t$ is $f = 1 - \exp(-k_1 t)$. At $t = 3600$~s (60~min):

\begin{equation}
  f = 1 - \exp(-1.7 \times 10^{-5} \times 3600)
    = 1 - e^{-0.061}
    \approx 5.9\%
  \label{eq:f_exchange}
\end{equation}

Equation~\eqref{eq:f_exchange} is a strict upper bound: Schlenoff et al.\ further show
that exchange in aqueous conditions is negligible, with essentially no exchange detected
even after 125~h in water (Figures~3--6 of \cit{7}). They conclude that desorption is
``least pronounced in aqueous systems.'' Wound fluid is an aqueous matrix; thiol exchange
does not contribute meaningfully to SAM loss over the assay window.

\subsection{Electrochemical desorption}

Reductive desorption of alkanethiol SAMs on gold occurs via:

\begin{equation}
  \text{Au--SR} + e^- \longrightarrow \text{Au}^0 + \text{RS}^-
  \label{eq:red_des}
\end{equation}

and is only possible when the electrode potential reaches or exceeds the desorption
potential $E_{\mathrm{des}}$. Widrig et al.\ measured $E_{\mathrm{des}}$ for a series of
$\omega$-alkanethiol SAMs on gold by cyclic voltammetry in 0.5~M KOH, finding a linear
dependence on carbon chain length $n$ \cit{6}:

\begin{equation}
  E_{\mathrm{des}} = (-0.83 - 0.020\,n)\ \text{V vs Ag/AgCl}
    \quad (0.5\ \text{M KOH})
  \label{eq:edes}
\end{equation}

For MCH ($n = 6$):

\begin{equation}
  E_{\mathrm{des}} = -0.83 - 0.020 \times 6 = -0.83 - 0.12 = -0.95\ \text{V}
  \label{eq:edes_mch}
\end{equation}

Love et al.\ report that desorption of alkanethiol SAMs from gold typically occurs in the
range $-1.0 \pm 0.25$~V vs Ag/AgCl \cit{1}, giving a conservative lower bound of
$-0.75$~V. The SWV scan window runs from $-0.65$~V to $+0.1$~V (Layer~6). The margin
between the scan minimum and the conservative lower bound is:

\begin{equation}
  \Delta E = E_{\mathrm{des,min}} - E_{\mathrm{scan,min}}
           = -0.75\ \text{V} - (-0.65\ \text{V})
           = -0.10\ \text{V}
  \label{eq:margin}
\end{equation}

The scan potential is 100~mV more positive than the most conservative desorption
estimate; the Widrig value for MCH is 300~mV more negative than the scan limit.
Electrochemical desorption does not occur during the assay.

The foregoing analysis uses Widrig's formula and Love's range, both derived for
simple $n$-alkanethiol SAMs whose hydrophobic chains have minimal solvent penetration
and desorb as a sharp stripping peak from a well-ordered layer. The probe-thiol
(HS--(CH$_2$)$_6$--PEG$_6$--SGKIPRTLTA--MB) and MCH have significant hydrophilic
character: the PEG$_6$ block is a polyether that binds water, and the peptide headgroup
bears partial charges at pH~7.4. A solvated, structurally disordered layer has different
desorption thermodynamics: the departing thiolate is stabilised by water, lowering
the thermodynamic barrier relative to a hydrophobic chain. The actual $E_{\text{des}}$
for the probe-thiol/MCH monolayer may therefore be less negative than the Widrig and
Love values indicate. This is an unconfirmed transferability issue; see Limitation~4.

The following assumptions underlie this section:

\begin{enumerate}[noitemsep]
  \item The oxidative stability argument uses the Schoenfisch O$_3$/air benchmark as an
    upper bound; no quantitative H$_2$O$_2$ rate constant is claimed (Limitation~2).
  \item Equation~\eqref{eq:k1} and the 5.9\% bound apply in organic solvent; aqueous
    exchange is negligible per Schlenoff \cit{7} (Limitation~3).
  \item The Love range \cit{1} is used as the primary reference for device-condition
    $E_{\mathrm{des}}$; the Widrig formula (equation~\ref{eq:edes}) is supporting data
    measured specifically in alkaline electrolyte (Limitation~4).
\end{enumerate}

%=======================================================================
\section{Ca$^{2+}$ Artefact and Ratiometric Cancellation}
\label{sec:calcium}

\subsection{Mechanism}

Wound fluid contains ionised Ca$^{2+}$ at a concentration comparable to serum
($\sim$1--2~mM; no paper in \texttt{papers/layer\_07/} reports a quantitative
wound-fluid Ca$^{2+}$ value: this is an unconfirmed estimate used for the sensitivity
calculation below).
If the probe-thiol architecture includes a PEG spacer between the alkyl chain and the
peptide, Ca$^{2+}$ can coordinate the polyether oxygen atoms in a crown-ether-like
interaction, contracting the spacer and decreasing the probe-to-gold bridge length by
$\Delta d$. Electron transfer (ET) from MB to gold through the bridge follows a
tunnelling mechanism, so a decrease in bridge length raises the apparent ET rate
constant $k_{\mathrm{ET}}$ and increases the SWV peak current of \emph{uncleaved}
probes.

For the peptide cleavage sensor (signal-off: MB released on cleavage), this
Ca$^{2+}$-induced increase in $k_{\mathrm{ET}}$ acts on the uncleaved probe population
equally before ($I_0$) and after ($I(t)$) the assay. The measured quantity is:

\begin{equation}
  \frac{\Delta I}{I_0} = \frac{I_0 - I(t)}{I_0}
  \label{eq:signal_off_ca}
\end{equation}

If Ca$^{2+}$ multiplies $k_{\mathrm{ET}}$ by a factor $R$ for all uncleaved probes
identically, then $I_0 \to R\,I_0$ and $I(t) \to R\,(1-f)\,I_0$ where $f$ is the
fraction cleaved. Substituting:

\begin{equation}
  \frac{\Delta I}{I_0} = \frac{R\,I_0 - R(1-f)\,I_0}{R\,I_0} = f
  \label{eq:ca_cancel}
\end{equation}

The Ca$^{2+}$ factor $R$ cancels exactly in the signal-off ratio, provided both $I_0$
and $I(t)$ are measured in the same buffer (same Ca$^{2+}$ concentration). The signal
is inherently self-normalising against Ca$^{2+}$-induced changes in electron transfer
rate, as long as the measurement protocol does not change the buffer between baseline
and post-incubation reads.

\subsection{Tunnelling decay constant}

The ET rate constant through a molecular bridge of length $d$ is:

\begin{equation}
  k_{\mathrm{ET}} = k_{\mathrm{ET}}^0\,\exp(-\beta\,d)
  \label{eq:tunnel}
\end{equation}

where $\beta$ is the tunnelling decay constant. Smalley et al.\ measured $\beta$ for
ferrocene-terminated (Fc) alkanethiol chains tethered to gold in 1.0~M HClO$_4$, by
fitting the chain-length dependence of the standard rate constant $k^\circ$ \cit{8}:

\begin{equation}
  \beta = 0.94 \pm 0.06\ \As
        = (9.4 \pm 0.6) \times 10^9\ \text{m}^{-1}
  \label{eq:beta}
\end{equation}

This measurement used ferrocene (Fc) as the redox probe, while the present device uses
methylene blue (MB). The value of $\beta$ is adopted here because $\beta$ is a property
of the alkyl bridge that electrons tunnel through, not of the terminal redox group.
Smalley et al.\ note that $\beta \approx 1.0\ \As$ is consistent across multiple redox
probes studied through the same alkyl bridge \cit{8}. The Fc-derived value therefore
transfers directly to the MB system.

\subsection{Sensitivity of $k_{\mathrm{ET}}$ to bridge contraction}

If Ca$^{2+}$ coordination contracts the PEG spacer by $\Delta d$, the fold increase in
$k_{\mathrm{ET}}$ follows from equation~\eqref{eq:tunnel}:

\begin{equation}
  \frac{k_{\mathrm{ET}}^{\mathrm{Ca}}}{k_{\mathrm{ET}}^0}
    = \exp(\beta\,\Delta d)
  \label{eq:kfold}
\end{equation}

For a representative contraction of $\Delta d = 0.2$~nm (one carbon--carbon bond
displacement, used as a sensitivity bound since $\Delta d$ is not measured for this
architecture; see Limitation~5):

\begin{align}
  \frac{k_{\mathrm{ET}}^{\mathrm{Ca}}}{k_{\mathrm{ET}}^0}
    &= \exp(9.4 \times 10^9\ \text{m}^{-1}
             \times 0.2 \times 10^{-9}\ \text{m}) \label{eq:kfold_plug} \\
    &= \exp(1.88) \approx 6.5 \notag
\end{align}

Both the active and control electrodes carry the same probe-thiol/MCH SAM and are
exposed simultaneously to the same wound fluid, so this fold increase applies identically
to both.

\subsection{Residual error from probe density change after cleavage}

Equation~\eqref{eq:ca_cancel} assumes that Ca$^{2+}$ affects uncleaved probes with the
same $k_{\mathrm{ET}}$ before and after the assay. This holds exactly if the cleaved
probe stub (the fragment remaining on the electrode after MB release) has negligible
electron transfer compared with the intact probe. If the stub retains some residual MB
(partial cleavage) or if cleavage changes the spacer geometry in a way that modifies
the Ca$^{2+}$ response of neighbouring probes, a second-order residual arises. The
magnitude of this residual depends on the specific peptide architecture and is not
quantifiable without a calibration measurement. It is expected to be small (second-order
in $f$ and in $\delta k_{\mathrm{ET,stub}}$), but experimental confirmation that
$\Delta I/I_0$ is Ca$^{2+}$-independent in the assay buffer is recommended.

\subsection{Limitations of this model}

\begin{enumerate}[noitemsep]

  \item \textbf{The deposition protocol and functional coverage limit are now confirmed;
    the absolute $\Gamma_p$ on the fabricated electrode must still be measured.}
    Wcis\l{}o~2021~\cit{11} provides a confirmed peptide-thiol deposition protocol
    (5~mM ethanolic solution, 20~h, MCH backfill). Adj\'{e}mian~2010~\cit{12} confirms
    the functional coverage regime ($\Gamma_0 \leq 10\ \text{pmol~cm}^{-2}$ for
    $>$80\% cleavage yield with MCH backfill). The device target of 2--5~pmol~cm$^{-2}$
    is within this regime. What remains unconfirmed is the \emph{actual} $\Gamma_p$
    achieved on the fabricated SGKIPRTLTA--MB electrode under the adopted deposition
    conditions. This must be measured directly by reductive desorption coulometry
    (method: Widrig~\cit{6}) or equivalent before Layer~8 and Layer~9 outputs can be
    quantified.

  \item \textbf{No quantitative $k_{\mathrm{ox}}$ for H$_2$O$_2$ oxidation of thiol SAMs.}
    Schoenfisch and Pemberton \cit{5} measure SAM oxidation by O$_3$ in air, not by
    H$_2$O$_2$ in solution. No paper in \texttt{papers/layer\_07/} provides $k_{\mathrm{ox}}$
    for aqueous H$_2$O$_2$. The oxidative stability argument is comparative only;
    experimental confirmation of SAM integrity after 60~min in wound-fluid H$_2$O$_2$ is
    required before the assay window can be certified.

  \item \textbf{Thiol exchange rate confirmed in hexane only.}
    Schlenoff et al.\ \cit{7} measure $k_1 = 1.7 \times 10^{-5}$~s$^{-1}$ in hexane
    and demonstrate negligible aqueous exchange qualitatively; a quantitative rate constant
    for exchange by Cys or GSH in phosphate-buffered wound fluid at pH~7.4 is not available.

  \item \textbf{$E_{\mathrm{des}}$ measured in 0.5~M KOH on hydrophobic alkanethiols;
    both conditions differ from this device.}
    The Widrig formula (equation~\ref{eq:edes}) was determined at pH~$\approx 13.7$ \cit{6}.
    At wound-fluid pH~7.4 the Nernst equation predicts that $E_{\mathrm{des}}$ shifts
    approximately $+(13.7 - 7.4) \times 0.059 \approx +0.37$~V toward positive values
    (because RS$^-$ is protonated at pH~7.4), placing $E_{\mathrm{des}}$ near
    $-0.95 + 0.37 = -0.58$~V. This is only 70~mV positive of the scan minimum of
    $-0.65$~V, compared with the 300~mV margin in alkaline electrolyte.
    The Love range \cit{1} ($-1.0 \pm 0.25$~V, conditions unspecified) is used as the
    primary working reference.
    A second applicability concern arises from molecular structure: Widrig's formula was
    measured on hydrophobic $n$-alkanethiols with no polar headgroups and no solvent
    penetration into the SAM interior. The probe-thiol and MCH in this device have
    substantial hydrophilic character (PEG$_6$ polyether, peptide headgroup, hydroxyl
    terminus of MCH). Greater water exposure and structural disorder lower the
    thermodynamic barrier to desorption and may shift $E_{\text{des}}$ toward less
    negative values relative to the Widrig formula. The formula is not transferable to
    this mixed peptide-thiol/MCH monolayer; experimental determination of
    $E_{\text{des}}$ in the assay buffer on the fabricated electrode is required before
    the scan window is certified safe.

  \item \textbf{Ca$^{2+}$ self-cancellation in $\Delta I/I_0$ is conditional on the
    measurement protocol.}
    Equation~\eqref{eq:ca_cancel} shows that the Ca$^{2+}$ artefact cancels in
    $\Delta I/I_0$ only if $I_0$ and $I(t)$ are measured in the same buffer (same
    Ca$^{2+}$ concentration). If the baseline $I_0$ is measured before adding wound
    fluid (in calibration buffer without Ca$^{2+}$) and $I(t)$ is measured in wound
    fluid (with Ca$^{2+}$), the cancellation fails and a spurious signal arises.
    Protocol design must ensure that both measurements are in the same ionic environment.
    Whether a PEG spacer is present and its Ca$^{2+}$ affinity are also unconfirmed for
    this probe architecture.

\end{enumerate}

%=======================================================================
\section*{References}

\noindent All papers in \texttt{electrochemical\_device/papers/layer\_07/}.

\begin{enumerate}[label={[\arabic*]},noitemsep,leftmargin=*]

\item Love JC, Estroff LA, Kriebel JK, Nuzzo RG, Whitesides GM (2005). Self-assembled
  monolayers of thiolates on metals as a form of nanotechnology.
  \textit{Chem Rev} 105(4):1103--1169.
  \textbf{love2005.pdf} -- $\Gamma_{\max}$ cross-check
  ($\approx 750$~pmol~cm$^{-2}$, p.~1108), $E_{\mathrm{des}} = -1.0 \pm 0.25$~V
  (p.~1122) confirmed from text.

\item Herne TM, Tarlov MJ (1997). Characterization of DNA probes immobilized on gold surfaces.
  \textit{J Am Chem Soc} 119(38):8916--8920.
  \textbf{herne1997.pdf} -- \textit{DNA-thiol specific. KH$_2$PO$_4$ buffer requirement,
  log-time deposition kinetics, and MCH displacement of non-thiolated DNA all pertain to
  HS-ssDNA with 25 charged phosphate groups per strand. These results do not transfer to
  the peptide-thiol probe and are not used for deposition protocol in this layer.}

\item Steel AB, Herne TM, Tarlov MJ (1998). Electrochemical quantitation of DNA immobilized
  on gold.
  \textit{Anal Chem} 70(22):4670--4677.
  \textbf{steel1998.pdf} -- \textit{DNA-thiol specific. The functional coverage limit
  ($4 \times 10^{12}$~molecules~cm$^{-2}$) was determined by hybridisation steric
  hindrance. Enzyme-cleavage steric limits for MMP-9 are different and not covered
  by this paper.}

\item Petrovykh DY, Kimura-Suda H, Whitman LJ, Tarlov MJ (2003). Quantitative analysis and
  characterization of DNA immobilized on gold.
  \textit{J Am Chem Soc} 125(17):5219--5226.
  \textbf{petrovykh2003.pdf} -- \textit{DNA-thiol specific. Log-linear adsorption kinetics
  and the non-Langmuir conclusion apply to HS-ssDNA. Peptide-thiol adsorption kinetics are
  not covered by this paper.}

\item Schoenfisch MH, Pemberton JE (1998). Air stability of alkanethiol self-assembled
  monolayers on silver and gold surfaces.
  \textit{J Am Chem Soc} 120(18):4502--4513.
  \textbf{schoenfisch1998.pdf} -- O$_3$ as primary oxidant confirmed from text; C$_{12}$
  SAM 36.9\% sulphur loss in 1~h confirmed from Table~5. Paper does not give $k_{\mathrm{ox}}$
  for aqueous H$_2$O$_2$; value attributed to this paper in \texttt{sam\_stability.py} is not
  confirmed and has been removed.

\item Widrig CA, Chung C, Porter MD (1991). The electrochemical desorption of
  $n$-alkanethiol monolayers from polycrystalline Au and Ag electrodes.
  \textit{J Electroanal Chem} 310(1--2):335--359.
  \textbf{widrig1991.pdf} -- $\Gamma_{\max} = 9.3 \times 10^{-10}$~mol~cm$^{-2}$
  (p.~351) confirmed; $E_{\mathrm{des}}$ formula (equation~\ref{eq:edes}) confirmed from
  text. Measured in 0.5~M KOH; pH correction noted in Limitation~4.

\item Schlenoff JB, Li M, Ly H (1995). Stability and self-exchange in alkanethiol
  monolayers.
  \textit{J Am Chem Soc} 117(50):12528--12536.
  \textbf{schlenoff1995.pdf} -- $k_1 = 1.7 \times 10^{-5}$~s$^{-1}$ in hexane confirmed
  from Table~3; concentration-independence of exchange rate confirmed from text; negligible
  aqueous exchange confirmed from Figures~3--6. Bimolecular rate model in
  \texttt{sam\_stability.py} is not supported by this paper and has been removed.

\item Smalley JF, Feldberg SW, Chidsey CED, Linford MR, Newton MD, Liu Y-P (2003). The
  kinetics of electron transfer through ferrocene-terminated alkanethiol monolayers on gold.
  \textit{J Am Chem Soc} 125(7):2004--2013.
  \textbf{smalley2003.pdf} -- $\beta = 0.94 \pm 0.06\ \As$ confirmed from text (p.~2009).
  Measured using ferrocene (Fc) redox probe; value transfers to methylene-blue system because
  $\beta$ is a property of the alkyl bridge, confirmed from discussion in \cit{8}.

\item Becka AM, Miller CJ (1993). Electrochemistry at $\omega$-hydroxy thiol coated
  electrodes. 3. Voltage independence of the electron tunneling barrier and measurements
  of redox kinetics at large overpotentials.
  \textit{J Phys Chem} 97(23):6233--6239.
  \textbf{becka1993.pdf} -- provides context for tunnelling through alkanethiol SAMs;
  the $\beta = 0.94\ \As$ value used in this layer is from Smalley~2003~\cit{8}, not
  Becka~1993.

\item Kridel SJ, Chen E, Kotra LP, Howard EW, Mobashery S, Smith JW (2001). Substrate
  hydrolysis by matrix metalloproteinase-9.
  \textit{J Biol Chem} 276(23):20572--20578. DOI:~10.1074/jbc.M100900200.
  \textbf{PIIS0021925819405115.pdf} -- $k_{\text{cat}} = 86.4\ \text{s}^{-1}$,
  $K_m = 1.3\ \text{mM}$, $k_{\text{cat}}/K_m = 67{,}500\ \Mms$ for peptide A11
  (SGKIPRTLTA) confirmed from Table~II. These solution-phase values are an upper bound
  for the surface $k_{\text{eff,surf}}$.

\item Wcis\l{}o A, Ma{\l}uch I, Niedzia{\l}kowski P, Ossowski T, Prahl A (2021).
  Label-free electrochemical test of protease interaction with a peptide substrate
  modified gold electrode.
  \textit{Chemosensors} 9(8):199. DOI:~10.3390/chemosensors9080199.
  \textbf{10.3390@chemosensors9080199.pdf} -- deposition protocol for peptide-thiol
  on gold (5~mM EtOH, 20~h, backfill with alkanethiol, BSA block) confirmed from
  Section~2.3.

\item Adj\'{e}mian J, Anne A, Cauet G, Demaille C (2010). Cleavage-sensing redox peptide
  monolayers for the rapid measurement of the proteolytic activity of trypsin and
  $\alpha$-thrombin enzymes.
  \textit{Langmuir} 26(12):10347--10356. DOI:~10.1021/la100397g.
  \textbf{10.1021@la100397g.pdf} -- cleavage yield $>$80\% for $\Gamma_0 \leq
  10\ \text{pmol~cm}^{-2}$ with MCH backfill confirmed from Fig.~2 and text (p.~10351);
  working coverage $(0.5\text{--}1) \times 10^{-11}$~mol~cm$^{-2}$ confirmed. Architecture
  (Fc-peptide-PEG tethered on Au, MCH backfill) is directly analogous to this device.

\item Sfragano PS, Moro G, Polo F, Palchetti I (2021). The role of peptides in the
  design of electrochemical biosensors for clinical diagnostics.
  \textit{Biosensors} 11(8):246. DOI:~10.3390/bios11080246.
  \textbf{sfragano2021.pdf} -- review of peptide-based electrochemical biosensors;
  confirms PEG linker strategies, cysteine-gold anchoring, and protease cleavage
  signal-off mechanisms consistent with this device architecture.

\end{enumerate}

\end{document}
